\section{Geometry of the precoloring}
\label{sec:geometry}

In this section we describe the \(15\)-coloring of \(\OG_6^*\)
constructed in Section~\ref{sec:chromatic} in geometric terms.  Keep
the notation of that section: \(T\) is the totally isotropic
\(3\)-space
\[
T=\Span\{e_1+e_2,\ e_3+e_4,\ e_5+e_6\},
\]
\(\mathcal C\) is the \(15\)-clique formed by the nonzero subspaces of
\(T\), and each \(W\in\mathcal C\) carries its own color \(c_W\).

For a vertex \(U\) of \(\OG_6^*\) outside \(\mathcal C\), define
\[
H(U)=U^\perp\cap T,\qquad K(U)=U\cap T.
\]
Since \(T^\perp=T\) and \(U\not\subseteq T\), we have \(H(U)\ne T\), so
\(\dim H(U)\in\{0,1,2\}\).  The following observation is the key to
the precoloring.

\begin{lemma}\label{lem:H-adjacency}
Let \(U\) be a vertex of \(\OG_6^*\) outside \(\mathcal C\) and let
\(W\in\mathcal C\).  Then \(U\) is adjacent to \(W\) if and only if
\(W\subseteq H(U)\).  Consequently, the allowed color list of \(U\)
is
\[
L(U)=\{c_W : 0\ne W\le T,\ W\not\subseteq H(U)\},
\]
and its size is
\[
|L(U)|=\begin{cases}
15,&\dim H(U)=0,\\
14,&\dim H(U)=1,\\
11,&\dim H(U)=2.
\end{cases}
\]
\end{lemma}

\begin{proof}
By definition, \(U\) is adjacent to \(W\) if and only if
\(W\subseteq U^\perp\).  Since \(W\subseteq T\), this is equivalent to
\(W\subseteq U^\perp\cap T=H(U)\).  The list size is
\(15-N(\dim H(U))\), where \(N(0)=0\), \(N(1)=1\), and \(N(2)=4\);
this gives \(15\), \(14\), and \(11\) in the three cases.
\end{proof}

The distribution of the \(2809=2824-15\) vertices outside
\(\mathcal C\), grouped by the triple
\((\dim U,\dim H(U),\dim K(U))\), is recorded in
Table~\ref{tab:profiles}.

\begin{table}[htbp]
\centering
\begin{tabular}{c|r|r|r}
\toprule
\(\dim U\) & \(\dim H(U)\) & \(\dim K(U)\) & Vertices \\
\midrule
\(1\) & \(2\) & \(0\) & \(56\) \\
\(2\) & \(1\) & \(0\) & \(448\) \\
\(2\) & \(2\) & \(1\) & \(196\) \\
\(3\) & \(0\) & \(0\) & \(512\) \\
\(3\) & \(1\) & \(1\) & \(784\) \\
\(3\) & \(2\) & \(2\) & \(98\) \\
\(4\) & \(0\) & \(1\) & \(448\) \\
\(4\) & \(1\) & \(2\) & \(196\) \\
\(4\) & \(2\) & \(3\) & \(7\) \\
\(5\) & \(0\) & \(2\) & \(56\) \\
\(5\) & \(1\) & \(3\) & \(7\) \\
\(6\) & \(0\) & \(3\) & \(1\) \\
\bottomrule
\end{tabular}
\caption{Distribution of the \(2809\) vertices outside the
\(15\)-clique, according to the triple
\((\dim U,\dim H(U),\dim K(U))\).}
\label{tab:profiles}
\end{table}

The \(357\) vertices with \(\dim H(U)=2\) are the most constrained,
having only \(11\) allowed colors.  A direct analysis of the
incidence structure between these vertices and the \(15\) clique
colors shows that the precoloring extends; this is the content of
the certificate described in Section~\ref{sec:chromatic}.

\subsection{Towards a structural description}

The invariants \((\dim U,\dim H(U),\dim K(U))\) do not by themselves
determine the color in the certificate.  Grouping the archived
coloring by this triple produces \(240\) distinct profiles outside
the clique; among these, \(139\) receive more than one color.  The
remaining \(101\) profiles receive a single color.  This is the
precise sense in which the present coloring depends on finer data
than the basic geometric invariants.

Two directions appear promising for a structural description.

\begin{enumerate}
\item \emph{Enrichment of the invariants.}  The profiles that receive
a single color suggest that the triple
\((\dim U,\dim H(U),\dim K(U))\) captures part of the coloring rule.
One may look for an additional invariant, such as the position of
\(K(U)\) inside \(H(U)\) or the structure of the quotient
\(U/(U\cap H(U))\), that separates the multi-color profiles.

\item \emph{The stabilizer of \(T\).}  The orthogonal group of
\(V_6\) acts on \(\OG_6^*\) and permutes the \(15\) clique colors.
A coloring invariant under the stabilizer of \(T\) in this group
would reduce the problem to a finite group-theoretic computation.
This direction has not been pursued in the present certificate.
\end{enumerate}

A complete structural proof that \(\chi(\OG_6^*)=15\), avoiding the
finite computation, remains open.  We record this as
Problem~\ref{prob:structural} in Section~\ref{sec:open}.
